\documentclass[10pt]{article}
\usepackage[margin=0.72in]{geometry}
\usepackage[T1]{fontenc}
\usepackage[utf8]{inputenc}
\usepackage{lmodern}
\usepackage[hidelinks]{hyperref}
\setlength{\parindent}{0pt}
\setlength{\parskip}{0.55em}
\pagestyle{empty}
\sloppy
\begin{document}
\begin{flushright}
Hamilton Jeong\\
College of Agriculture and Life Sciences, Cornell University\\
215 Garden Avenue, Ithaca, NY 14850, USA\\
\href{mailto:hhj22@cornell.edu}{hhj22@cornell.edu}\\
16 October 2026
\end{flushright}

Editors-in-Chief\\
\emph{International Journal of Computer Assisted Radiology and Surgery}

Dear Editors-in-Chief,

Please consider our Original Article, ``Operator identity and out-of-user generalization in automated surgical skill assessment: a multi-dataset evaluation,'' for publication in the \emph{International Journal of Computer Assisted Radiology and Surgery}.

Automated surgical-performance studies often contain repeated trials from the same operators. This design creates an opportunity for models to exploit operator-specific motion patterns, particularly when development and evaluation records overlap by person. We evaluated this problem across 2,501 trials from four surgical-training datasets using paired same-participant and held-out-participant protocols, formal label permutations, fresh identity probes, prespecified identity-suppression interventions, negative controls, and participant-level inference.

The principal finding is deliberately narrow and actionable: operator identity was recoverable across all examined task groups, and targeted linear projection reduced identity recovery under fresh linear and nonlinear probes, but this verified reduction did not consistently improve prediction for unseen operators. The work therefore shows why identity recoverability, successful suppression, and held-out-user utility must be evaluated as separate claims. The study aligns directly with IJCARS interests in surgical simulation and training, machine intelligence, and methods of validation and verification.

The manuscript transparently discloses overlap with our earlier JIGSAWS-only unpublished manuscript and software repository, which analyzed the same 103 JIGSAWS trials. The present article is a replication and mechanistic extension rather than independent JIGSAWS validation: it adds three external datasets, a harmonized 41-variable pipeline, direct representation interventions, fresh post-intervention probes, matched controls, and participant-bootstrap inference. It does not reuse the earlier commercial-model, calibration, or selective-prediction analyses.

The manuscript is original, is not under consideration elsewhere, and has been approved by all authors. The authors received no funding for this work and declare no financial or non-financial competing interests. This was a secondary analysis of existing de-identified datasets; the investigators had no direct identifiers, re-identification keys, or participant contact. Under Cornell University guidance, the activity did not constitute federally regulated human-participant research and did not require further Cornell IRB review. Code, configurations, aggregate results, figures, and artifact hashes are available in the linked GitHub repository, subject to the source datasets' access terms.

Thank you for your consideration.

Sincerely,\\[1.5em]
Hamilton Jeong\\
Corresponding author
\end{document}
